\documentclass[11pt]{article}

\usepackage[margin=1in]{geometry}
\usepackage{amsmath}
\usepackage{booktabs}
\usepackage{array}

\title{Motivation Note for the Current InfoGate Framework}
\author{}
\date{April 11, 2026}

\begin{document}

\maketitle

\section{Introduction}

Multimodal sentiment analysis aims to predict sentiment from language,
acoustic, and visual signals. Compared with unimodal sentiment analysis, it can
capture richer affective evidence, but it also introduces a harder fusion
problem: different modalities do not contribute equally, their sequence quality
is highly uneven, and their usefulness can change from one sample to another.
In practice, the main difficulty is not simply how to combine more modalities,
but how to identify which modality should dominate the current decision and how
to prevent noisy auxiliary information from corrupting that dominant stream.

This issue is especially visible in primary-centric multimodal sentiment
methods. Prior work has shown that treating all modalities symmetrically can
blur the most informative signal, while fixed language-centric strategies are
too rigid for samples where acoustic or visual cues carry the decisive emotion.
Dynamic primary-modality selection is therefore a promising direction.
However, once the primary modality is allowed to vary across samples, a new
problem appears: the selector and fusion module can only be as reliable as the
features they consume. If those features still contain task-irrelevant
variation, sequential redundancy, or uncertain cross-modal evidence, then both
primary-modality selection and cross-modal injection become unstable.

The current InfoGate framework is motivated by this gap. Instead of asking only
which modality should be primary, it asks a stricter question: how can we
first purify multimodal representations into a more task-relevant space, then
perform dynamic primary selection inside that space, and finally inject
auxiliary information into the primary stream in a reliability-aware manner?
Under this view, the core problem is not ordinary multimodal fusion, but
\emph{noise-aware, sample-adaptive, primary-centric fusion}. The framework is
designed to solve exactly this problem.

\section{Related Work}

\noindent \textbf{Symmetric multimodal fusion.} Many multimodal sentiment methods
focus on stronger cross-modal interaction while treating the three modalities in
a roughly symmetric way. This line has produced effective fusion architectures,
but it implicitly assumes that all modalities deserve comparable trust for the
current sample. When one modality is much cleaner or more emotionally decisive
than the others, symmetric fusion can mix useful evidence with distraction.

\noindent \textbf{Static language-centric modeling.} Another line of work gives
language a persistent dominant role, since text often has the highest semantic
density in sentiment analysis. This is a reasonable prior, but it is still a
global prior. It breaks down on samples where sentiment is expressed mainly by
voice or facial dynamics rather than literal wording.

\noindent \textbf{Dynamic primary-modality selection.} More recent work, such as
MODS, moves beyond fixed modality priority and argues that the dominant modality
should be selected dynamically for each sample. This is an important step, but
it still leaves two open questions: what kind of features should be used for
selection, and how should auxiliary modalities be injected after the primary
stream is chosen? If raw or weakly filtered representations are used, then
dynamic selection alone does not fully solve the noise problem.

\noindent \textbf{Information bottleneck and informative spaces.} ITHP and CyIN
suggest a complementary direction: rather than trusting raw multimodal features,
the model should compress them into a more informative latent space that keeps
task-relevant content and suppresses nuisance variation. This idea is especially
appealing for multimodal sentiment analysis, where acoustic and visual streams
often contain long, redundant, and unstable sequences. The current InfoGate
framework follows this intuition and combines it with primary-centric fusion.

\section{Motivation of the Current Framework}

The current framework is motivated by three concrete observations.

First, sample-wise modality dominance is real, but dynamic selection is only
useful if the selector sees representations that are already reasonably clean.
If the acoustic or visual stream remains cluttered with redundancy, or if the
text stream contains sentiment-irrelevant lexical variation, then the routing
decision can be unstable. A primary selector should therefore act on
task-oriented features rather than on raw projected features.

Second, auxiliary modalities are not uniformly helpful even after a primary
stream has been selected. Standard cross-attention assumes that all auxiliary
tokens are equally eligible to influence the primary representation, which is too
strong an assumption for multimodal sentiment data. Some auxiliary tokens are
supportive, some are neutral, and some are actively harmful. The fusion module
therefore needs a way to reflect uncertainty and to attenuate unreliable
cross-modal evidence instead of injecting it indiscriminately.

Third, if the final prediction is meant to reflect the dominant sentiment signal,
then the last-stage representation should remain primary-centric. Directly
concatenating large residual features from text or from all modalities can make
the classifier succeed by bypassing the intended routing logic, which weakens
the meaning of primary-modality modeling. A cleaner design is to let the model
predict from the enhanced primary bottleneck stream itself.

These observations motivate the present framework: an information-filtered,
confidence-aware, dynamically routed, primary-centric multimodal sentiment
architecture.

\section{What Problem Is Solved, and What Solution Is Proposed}

The framework can be summarized as a problem--solution mapping shown in Table~1.

\begin{table}[h]
\centering
\begin{tabular}{>{\raggedright\arraybackslash}p{0.30\linewidth} >{\raggedright\arraybackslash}p{0.60\linewidth}}
\toprule
Problem & Proposed solution in the current framework \\
\midrule
Sample-dependent modality dominance & Use MSelector to choose the primary modality dynamically for each sample instead of fixing language as the default center. \\
Noisy and task-irrelevant unimodal features & Apply Information Bottleneck encoders to each modality so that routing and fusion operate on compressed latent representations rather than raw projected features. \\
Auxiliary modalities have uneven reliability & Derive per-token confidence from the bottleneck uncertainty and use it to modulate cross-attention, so uncertain auxiliary evidence is down-weighted. \\
Auxiliary information should not be injected with equal strength & Introduce adaptive information gates that control how much each auxiliary cross-attention branch contributes to the primary stream. \\
The final classifier may bypass primary-centric fusion & Remove the direct multimodal highway and predict only from the enhanced primary bottleneck representation. \\
Cross-modal semantics can drift apart during training & Use bottleneck reconstruction and CRA-style translation regularization to keep modality interaction in a more consistent informative space. \\
\bottomrule
\end{tabular}
\caption{Problem--solution view of the current InfoGate framework.}
\end{table}

Operationally, the framework first encodes text, acoustic, and visual inputs into
a shared hidden space. Each modality is then passed through an Information
Bottleneck encoder that produces both a compressed latent feature and a
confidence estimate. These bottlenecked features are pooled and sent to the
dynamic MSelector, which predicts a sample-specific routing distribution over
modalities and determines the primary stream. After routing, the two auxiliary
streams interact with the primary stream through confidence-modulated
cross-attention. Their contributions are further filtered by adaptive gates,
allowing the model to preserve helpful evidence while suppressing noisy or
misaligned information. Finally, the enhanced primary stream alone is aggregated
and passed to the prediction head, keeping the decision process consistent with
the primary-centric design.

\section{Discussion}

This motivation is broadly defensible. The problem statement is concrete, the
modules correspond to identifiable failure modes, and the overall framework has
a clear internal logic: first filter information, then select the dominant
stream, then fuse auxiliary evidence under uncertainty control, and finally make
the decision from the enhanced primary representation. This is a stronger story
than merely saying the model uses several good modules together.

At the same time, the framing should remain disciplined. The current branch is
best presented as a solution for complete-modality multimodal sentiment
analysis, with dynamic primary selection and noise-aware fusion as the central
contributions. Although the code still contains translation-style regularization,
the present version should not be overclaimed as a full missing-modality system
unless that setting is explicitly evaluated again. Likewise, the motivation is
strongest when it emphasizes \emph{task-relevant information control} and
\emph{primary-centric fusion stability}, rather than making broad claims about
all multimodal learning settings.

In short, the motivation can stand. A clean one-sentence formulation is the
following: existing primary-centric MSA methods recognize that different samples
need different dominant modalities, but they do not sufficiently control the
quality and reliability of the information entering selection and fusion;
therefore, the current framework proposes an Information Bottleneck-guided,
confidence-aware, dynamically routed primary-centric architecture to make
modality selection and cross-modal interaction more stable and task-relevant.

\end{document}